\section{Conclusion}

From a theoretical point of view, holographic displays fulfil the requirements of the ultimate display outlined in section \ref{sec:definition} perfectly. In practice, the amount of data contained in such an uncompromising hologram will likely remain impossible to generate, transmit and display in real time for decades to come. Super-multiview displays also have a claim as they are effectively an approximation of holography using 2D imaging techniques. But as such, they suffer from the same challenges regarding data-rate as holograms. Volumetric displays provide all physiological depth cues, but fail to reach voxel densities competitive with 2D displays and struggle with occlusion and texture. They also don't fit the (admittedly arbitrary) concept of a \textit{window view upon the world}. Out of all the displays presented so far, the one presented by An et al. \cite{An_Won_Kim_Hong_Kim_Kim_Song_Choi_Kim_Seo_2020} comes closest. It respects all parameters of the plenoptic function and provides every depth cue except for vertical parallax. If innovations in computational display holography continue, it seems plausible that the differences to an uncompromising hologram become negligible within ten years.